\section{Introducción}

\begin{frame}{Autómata celular}
	\begin{center}
		\begin{tikzpicture}
			\node[cajaobjetivo, text width=0.85\textwidth]
			{
				\begin{definition}
					
					Un autómata celular es una 4-tupla ($\mathcal{L}, \mathcal{S}, \mathcal{N}, \mathcal{R}$) donde: \\
					
					\begin{itemize}
						
						\item $\mathcal{L} \subseteq \mathbb{Z}^D$ es el espacio celular $D$-dimensional. 
						
						\item $\mathcal{S}$ es un conjunto de estados. 
						
						\item $\mathcal{N} = (\vec{v}_1, \vec{v}_2, \dots, \vec{v}_m )$ es el conjunto de vectores de desplazamiento que definen la vecindad de una célula en $\mathcal{L}$. 
						
						\item $\mathcal{R} : \mathcal{S}^m \to \mathcal{S}$ es la regla local del autómata
						
						
					\end{itemize}
					
				\end{definition}
			};
		\end{tikzpicture}
	\end{center}
\end{frame}

\begin{frame}{Sistema complejo}
	\begin{center}
		\begin{tikzpicture}
			\node[cajaobjetivo, text width=0.85\textwidth]
			{
				\begin{definition}
					
					Un sistema complejo es una entidad compuesta de partes interconectadas, de tal forma que el comportamiento colectivo es más que la suma de sus componentes individuales.
					
				\end{definition}
			};
		\end{tikzpicture}
	\end{center}
\end{frame}
